\section{Tencentが視覚・PC操作対応のHy4を発表}
Tencentはフラッグシップモデル「Hy4」を正式発表した。Grokの収集によれば、7700億パラメータ規模、最大100万トークンのコンテキストを備え、画像理解に加えてCLIやGUIを操作する機能を持つ。金融、フロントエンド開発、ゲーム、研究などへの応用を訴求し、WorkBuddyやCodeBuddyでの期間限定無料提供にも言及した。オープンウェイト公開は「近日中」とされ、現時点で公開済みとは扱えない。早期利用者の改善評価や性能比較はあるものの、独立ベンチマーク、ライセンス、安全対策の詳細は今後の確認事項である。
\dailyxsource{Tencent Hy、11月4日15:43 JST}{https://x.com/TencentHunyuan/status/2117870250838630539}

\section{Claudeがvault粒子の構造生物学研究を支援}
Anthropicは、Claudeを活用して細菌由来のvault particleの三次元構造をマッピングした研究を紹介した。投稿に添えられたプレプリントでは、関連タンパク質の構造解明を通じ、従来は数週間から数か月を要した作業を数日へ短縮したと報告している。一方、Claudeが単純な誤りを犯す場面もあり、研究者との協働が必要だったことを明記した。AIによる科学研究の加速を示す事例として注目されるが、査読状況や実験データとの対応関係、独立した再現性は別途検討を要する。
\dailyxsource{Anthropic、11月5日05:11 JST}{https://x.com/AnthropicAI/status/2118073665296228863}

\section{Metaがエージェント管理ツールConductorを買収}
AIコーディングスタートアップConductorのCEO、Charlie Holtzが、同社のMetaへの参加をXで発表した。Conductorは複数のコーディングエージェントを並列実行するオーケストレーションツールを提供してきた。報道によれば、チームはMeta Superintelligence Labsで次世代インターフェース開発に取り組み、Conductor Cloudは11月20日に終了する一方、ローカルアプリは継続するという。買収額は非公開。開発者が多数のエージェントを扱うための操作基盤が、モデル自体と並ぶ競争領域になりつつある。
\dailyxsource{Charlie Holtz、11月5日03:18 JST}{https://x.com/charlieholtz/status/2118045232218743264}

\section{Mistral Large 4の脆弱性対応デモ}
MistralはLarge 4（Le Chonk）を用い、DjangoのCVEに対する脆弱性発見、実証、修正までを133のエージェントで進めるデモ動画を公開した。トリアージからパッチ作成までの一連の工程を短時間で示し、セキュリティ実務へのエージェント応用を訴えた。ただしGrokが確認した詳細資料は投稿内の動画が中心で、技術ブログや実行条件は未確認である。デモで示された作業がどこまで再現可能か、実際の脆弱性修正の品質や他モデルとの比較は今後の評価対象となる。
\dailyxsource{Mistral AI、11月4日19:24 JST}{https://x.com/MistralAI/status/2117925908572991748}

\section{LongCat-2.5など中国モデルの比較が続く}
MeituanのLongCat-2.5を含む中国モデル群について、観測窓内に比較・体験投稿が複数見られた。Grokの記録ではLongCat-2.5は1.6兆パラメータ規模のマルチモーダル・エージェントモデルと紹介され、SVG生成など具体的なタスクを使った性能比較や価格の議論が共有された。ただし、公式発表そのものは今回の観測窓外であり、期間内の話題は二次的な言及が中心である。個別利用者の優位性評価を一般的な性能差とみなすことはできず、公式指標と独立評価の照合が待たれる。
\dailyxsource{技術者による比較投稿（例）、観測窓内}{https://x.com/noclipepe/status/2118102232780480804}
